% Pursuing Stacks (cotes 134-2 to 134-8): archive cote 134-7, site page 106 (PDF page 107), typescript page 548.
% typescript. Draft compared: ps6.tex:5323-5408.
% AI-transcribed by Claude Opus 5.5 (claude-opus-5-5, adaptive thinking, effort high), Message Batches API, batch ps4 (msgbatch_012Aoc1NuskATPsyBeKSYK7J), 2026-11-12; not yet checked by eye.
\scanpage{134-7/106}
\mspage{548}

be the corresponding augmented k-algebra admitting M as a square-zero
augmentation ideal,  and let
\[ (7) \qquad I(M) = \mathrm{Spec}(D(M)) \]
it's spectrum,  which is a first order infinitesimal neighbourhood of
S = Spec(k).  For variable M,  D(M) depends covariantly, I(M) contravari-%
antly on M.  Thus
\[ (8) \qquad \xout{D(L^*)} \quad D(L^*) \]
is a co-semisimplicial algebra,  depending covariantly on $L^*$,  hence
contravariantly on $L_* = V(L^*)$,  and accordingly
\[ (9) \qquad I(L^*) = \mathrm{Spec}(D(L^*)) \]
(where in both members I and Spec are applied again componentwise) is
a semisimplicial affine scheme,  indeed a first order infinitesimal one,
depending contravariantly on $L^*$,  hence covariantly on $L_*$ .  Now the
last member of (*)  may be interpreted \xout{as} non-linearly as
\[ (**) \qquad \mathrm{Hom}\xout{Hom}_{\text{co-ss k-alg}}(A^*,D(L^*)) \hand{\simeq} \mathrm{Hom}(I(L^*), X_*) \quad , \]
where in the second term the Hom means homomorphisms of semisimplicial
schemes.  Let's now write (with an obvious afterthought !)
\[ (10) \qquad \hand{\widetilde{\phantom{S}}}\!\!\!S(L_*) \;=\; I(L^*) \qquad \text{whenever} \quad L_* \;=\; V(L^*) \quad , \]
the sequence of isomorphisms (*) and (**) may be summed up by
\[ (11) \qquad \mathrm{Hom}(L_*,L\unsure{\hand{\Pi}}_\bullet(X_*)) \hand{\simeq} \mathrm{Hom}(\hand{\widetilde{\phantom{S}}}\!\!\!S(L_*), X_*) \quad , \]
where the first Hom \deleted{is taken in Hotab$_o$(k)  i.e. it} \inserted[typed]{means} \xout{xxxx} maps of semisimpl%
\hand{-}cial $\underline{O}_k$-Modules,  whereas the second is a Hom of semisimplicial schemes.
It looks very much like the adjunction formula (p. 524,  (7 b)) we are
after,  with however two big grains of salt.  The smaller one is that
this formula does'nt take place in derived categories,  but rather in
the categories of would-be models.  The considerably bigger grain of
salt,  is that  $\hand{\widetilde{\phantom{S}}}\!\!\!S(L_*)$  is'nt at all in our model-category,  it's compo-%
nents are purely infinitesimal,  first order schemes,  and far from
being vector bundles !  Essentially,  what we have been using for get-%
ting the (admittedly quite tautological) \deleted{adjunction} \inserted{adjunction} formula (11),  is
that the L\unsure{\hand{i}}e functor on schemes-with-section  over S is representable
in an obvious way,  namely by the scheme-with-section
\[ I(k) = \mathrm{Spec}(k\hand{[}T\hand{]}/(T^2)) \quad , \]
which is a first\unsure{\hand{-}}order infinitesimal scheme-with-section over k.

\margin{106}
